% =====================================================================
% Two-column layout summarising the Jexpresso (Julia) -- Alya (Fortran)
% coupling. Each "row" pairs the matching snippets from the two codes
% that together implement one step of the coupling protocol. Snippets
% are lifted verbatim (lightly trimmed) from:
%   - AlyaProxy/myAlya.f90
%   - src/run.jl
%   - src/kernel/coupling/couplingStructs.jl
%
% Implementation note: lstlisting does NOT work inside tabular/tabularx
% p{} cells, because the cell is collected as a \parbox macro argument
% and verbatim-mode environments cannot be tokenised inside macro
% arguments. We therefore lay out the "table" with side-by-side
% minipages, whose bodies are processed normally and so accept
% lstlisting safely.
%
% Compile with: pdflatex coupling_snippets.tex
% =====================================================================
\documentclass[10pt,landscape]{article}

\usepackage[a4paper,landscape,margin=1.2cm]{geometry}
\usepackage{listings}
\usepackage{xcolor}
\usepackage{booktabs}

\definecolor{kw}{rgb}{0.10,0.30,0.65}
\definecolor{cm}{rgb}{0.35,0.55,0.35}
\definecolor{st}{rgb}{0.60,0.20,0.20}
\definecolor{bgF}{rgb}{0.97,0.97,0.92}
\definecolor{bgJ}{rgb}{0.93,0.96,0.99}
\definecolor{ruleC}{rgb}{0.55,0.55,0.55}

\lstdefinestyle{fortran}{
    language=Fortran,
    basicstyle=\ttfamily\scriptsize,
    keywordstyle=\color{kw}\bfseries,
    commentstyle=\color{cm}\itshape,
    stringstyle=\color{st},
    backgroundcolor=\color{bgF},
    breaklines=true,
    columns=fullflexible,
    keepspaces=true,
    showstringspaces=false,
    frame=single,
    framesep=2pt,
    aboveskip=2pt, belowskip=2pt,
}
\lstdefinelanguage{Julia}{
    morekeywords={function,end,if,else,elseif,for,while,return,do,in,using,import,
                  module,struct,const,let,global,local,true,false,nothing},
    morecomment=[l]{\#},
    morestring=[b]",
    sensitive=true,
}
\lstdefinestyle{julia}{
    language=Julia,
    basicstyle=\ttfamily\scriptsize,
    keywordstyle=\color{kw}\bfseries,
    commentstyle=\color{cm}\itshape,
    stringstyle=\color{st},
    backgroundcolor=\color{bgJ},
    breaklines=true,
    columns=fullflexible,
    keepspaces=true,
    showstringspaces=false,
    frame=single,
    framesep=2pt,
    aboveskip=2pt, belowskip=2pt,
}

% ------------------------------------------------------------------
% Layout helpers
% ------------------------------------------------------------------
% Total text width in landscape A4 with 1.2cm margins ~ 24.7cm.
% Layout: step label (3.0cm) + two code columns (~10.6cm each) +
% gutters.
\newlength{\stepW}\setlength{\stepW}{3.0cm}
\newlength{\codeW}\setlength{\codeW}{10.6cm}

% Row header (above each pair of code minipages).
\newcommand{\stepHdr}[1]{%
  \par\medskip
  \noindent\textcolor{ruleC}{\rule{\linewidth}{0.3pt}}\par
  \noindent\textbf{#1}\par\smallskip
}

\setlength{\parindent}{0pt}

\begin{document}

\begin{center}
{\Large\bfseries Jexpresso\,$\leftrightarrow$\,Alya coupling: matching code snippets}\\[2pt]
{\small Left column: \texttt{AlyaProxy/myAlya.f90}\quad
 Right column: \texttt{src/run.jl} and \texttt{src/kernel/coupling/couplingStructs.jl}}
\end{center}

\medskip

% Column headers
\noindent
\begin{minipage}[t]{\stepW}\centering\textbf{Step}\end{minipage}\hfill
\begin{minipage}[t]{\codeW}\centering\textbf{Alya (Fortran) --- \texttt{myAlya.f90}}\end{minipage}\hfill
\begin{minipage}[t]{\codeW}\centering\textbf{Jexpresso (Julia)}\end{minipage}

% ---------------------------------------------------------------- 1
\stepHdr{1.\ MPI init \& communicator split}
\noindent
\begin{minipage}[t]{\stepW}
\footnotesize Each code joins \texttt{MPI\_COMM\_WORLD} and creates a private communicator by colour.
\end{minipage}\hfill
\begin{minipage}[t]{\codeW}
\begin{lstlisting}[style=fortran]
call MPI_Init(ierr)
call MPI_Comm_rank(MPI_COMM_WORLD, rank, ierr)
call MPI_Comm_size(MPI_COMM_WORLD, size, ierr)

! Alya colour = 1
call MPI_Comm_split(MPI_COMM_WORLD, 1_4, rank, &
                    PAR_COMM_FINAL, ierr)
call MPI_Comm_size(PAR_COMM_FINAL, asize, ierr)
call MPI_Comm_rank(PAR_COMM_FINAL, arank, ierr)
\end{lstlisting}
\end{minipage}\hfill
\begin{minipage}[t]{\codeW}
\begin{lstlisting}[style=julia]
# run.jl : Jexpresso runs in MPMD with Alya, sharing
# MPI_COMM_WORLD.  A separate colour gives a local
# communicator for the Julia ranks only.
_world      = MPI.COMM_WORLD
_local_comm = MPI.Comm_split(_world, appid, wrank)
nparts      = MPI.Comm_size(_local_comm)
\end{lstlisting}
\end{minipage}

% ---------------------------------------------------------------- 2
\stepHdr{2.\ Handshake: application names}
\noindent
\begin{minipage}[t]{\stepW}
\footnotesize All ranks gather a 128-char name to world rank 0.  Jexpresso decides it is coupled iff \texttt{wsize>nparts}.
\end{minipage}\hfill
\begin{minipage}[t]{\codeW}
\begin{lstlisting}[style=fortran]
app_name = 'ALYA'
if (rank == 0) allocate(app_dumm(size))
call MPI_Gather(app_name, 128, MPI_CHARACTER, &
                app_dumm,  128, MPI_CHARACTER, &
                0, MPI_COMM_WORLD, ierr)
! rank 0 prints the world-rank -> code map
\end{lstlisting}
\end{minipage}\hfill
\begin{minipage}[t]{\codeW}
\begin{lstlisting}[style=julia]
# couplingStructs.jl :: je_perform_coupling_handshake
local_chars = Vector{UInt8}(rpad("JEXPRESSO", 128, ' '))
MPI.Gather!(local_chars, nothing, 0, world)
# coupling is active iff wsize > nparts
return MPI.Comm_size(world) > nparts
\end{lstlisting}
\end{minipage}

% ---------------------------------------------------------------- 3
\stepHdr{3.\ Grid metadata broadcast (Alya $\rightarrow$ Julia)}
\noindent
\begin{minipage}[t]{\stepW}
\footnotesize Alya rank 0 broadcasts the structured-grid bounding box and node counts in a fixed order; Jexpresso receives them in the same order.
\end{minipage}\hfill
\begin{minipage}[t]{\codeW}
\begin{lstlisting}[style=fortran]
rem_min = [-5000.0,-3000.0, 0.0]
rem_max = [ 5000.0, 1500.0, 10000.0]
rem_nx  = [10, 10, 10];  ndime = 3

call MPI_Bcast(ndime, 1, MPI_INTEGER, 0, &
               MPI_COMM_WORLD, ierr)
do idime = 1, 3
   call MPI_Bcast(rem_min(idime),1, &
        MPI_DOUBLE_PRECISION,0,MPI_COMM_WORLD,ierr)
   call MPI_Bcast(rem_max(idime),1, &
        MPI_DOUBLE_PRECISION,0,MPI_COMM_WORLD,ierr)
   call MPI_Bcast(rem_nx(idime), 1, &
        MPI_INTEGER,         0,MPI_COMM_WORLD,ierr)
end do
\end{lstlisting}
\end{minipage}\hfill
\begin{minipage}[t]{\codeW}
\begin{lstlisting}[style=julia]
# couplingStructs.jl :: je_receive_alya_data
ndime_buf = Vector{Int32}(undef, 1)
MPI.Bcast!(ndime_buf, 0, world)
ndime = Int(ndime_buf[1])

rem_min = Vector{Float64}(undef, 3)
rem_max = Vector{Float64}(undef, 3)
rem_nx  = Vector{Int32}(undef, 3)
for idime in 1:3
    MPI.Bcast!(@view(rem_min[idime:idime]), 0, world)
    MPI.Bcast!(@view(rem_max[idime:idime]), 0, world)
    MPI.Bcast!(@view(rem_nx[idime:idime]),  0, world)
end
\end{lstlisting}
\end{minipage}

% ---------------------------------------------------------------- 4
\stepHdr{4.\ Alya-to-world rank map (Allreduce)}
\noindent
\begin{minipage}[t]{\stepW}
\footnotesize Alya ranks publish their world rank; everyone (including Julia) reconstructs the map via \texttt{MPI\_Allreduce(SUM)}.
\end{minipage}\hfill
\begin{minipage}[t]{\codeW}
\begin{lstlisting}[style=fortran]
allocate(alya_to_world_snd(0:asize-1))
allocate(alya_to_world    (0:asize-1))
alya_to_world_snd        = 0
alya_to_world_snd(arank) = rank
call MPI_AllReduce(alya_to_world_snd, alya_to_world, &
     asize, MPI_INTEGER4, MPI_SUM, &
     MPI_COMM_WORLD, ierr)
\end{lstlisting}
\end{minipage}\hfill
\begin{minipage}[t]{\codeW}
\begin{lstlisting}[style=julia]
# Julia ranks contribute zeros; Allreduce on
# MPI_COMM_WORLD reconstructs the same map.
nranks_alya  = wsize - nparts
alya2world_l = zeros(Int32, nranks_alya)
alya2world   = MPI.Allreduce(alya2world_l,
                             MPI.SUM, world)
\end{lstlisting}
\end{minipage}

% ---------------------------------------------------------------- 5
\stepHdr{5.\ Point-count exchange (MPI\_Alltoall)}
\noindent
\begin{minipage}[t]{\stepW}
\footnotesize Julia tells every Alya rank how many of its grid points fall inside the Julia subdomain; Alya supplies zeros (it never initiates).
\end{minipage}\hfill
\begin{minipage}[t]{\codeW}
\begin{lstlisting}[style=fortran]
allocate(npoin_send(0:size-1))
allocate(npoin_recv(0:size-1))
npoin_send = 0   ! Alya never initiates
npoin_recv = 0

call MPI_Barrier(MPI_COMM_WORLD, ierr)
call MPI_Alltoall(npoin_send, 1, MPI_INTEGER4, &
                  npoin_recv, 1, MPI_INTEGER4, &
                  MPI_COMM_WORLD, ierr)
total_pts = sum(npoin_recv)
\end{lstlisting}
\end{minipage}\hfill
\begin{minipage}[t]{\codeW}
\begin{lstlisting}[style=julia]
# couplingStructs.jl : setup_coupling_and_mesh
# Each Julia rank counts how many Alya grid points
# fall inside its SEM subdomain, grouped by owner.
npoin_send = zeros(Int32, wsize)
for wr in alya_owner_ranks
    npoin_send[wr+1] += 1
end
send_to_ranks = Int32[i-1 for i in 1:wsize
                          if npoin_send[i] > 0]

recv_counts_from_alya = zeros(Int32, wsize)
MPI.Barrier(world)
MPI.Alltoall!(Vector{Int32}(npoin_send),
              recv_counts_from_alya, 1, world)
\end{lstlisting}
\end{minipage}

% ---------------------------------------------------------------- 6
\stepHdr{6.\ Send/Recv of Jexpresso SEM node global IDs}
\noindent
\begin{minipage}[t]{\stepW}
\footnotesize One-off transfer of the global indices of every Alya point each Julia rank will later send values for.  Alya stores them in \texttt{je\_gids\_all} and uses them to scatter into its local slot.
\end{minipage}\hfill
\begin{minipage}[t]{\codeW}
\begin{lstlisting}[style=fortran]
do i = asize, size-1
   if (npoin_recv(i) > 0) then
      call MPI_Recv(je_gids_all(je_gid_offset+1), &
           npoin_recv(i), MPI_INTEGER4, i, 0, &
           MPI_COMM_WORLD, MPI_STATUS_IGNORE, ierr)
      je_gid_offset = je_gid_offset + npoin_recv(i)
   end if
end do
\end{lstlisting}
\end{minipage}\hfill
\begin{minipage}[t]{\codeW}
\begin{lstlisting}[style=julia]
# couplingStructs.jl :: je_send_node_list
for dest_rank in send_to_ranks
    mask    = alya_owner_ranks .== dest_rank
    gid_buf = Int32.(alya_local_ids[mask])
    push!(send_requests,
          MPI.Isend(gid_buf, dest_rank, 0, world))
end
MPI.Waitall(send_requests)
\end{lstlisting}
\end{minipage}

% ---------------------------------------------------------------- 7
\stepHdr{7.\ Time loop: field exchange (Julia $\rightarrow$ Alya)}
\noindent
\begin{minipage}[t]{\stepW}
\footnotesize At every Jexpresso step, the SEM solution is interpolated to the Alya points and sent; Alya posts matching non-blocking receives and waits.
\end{minipage}\hfill
\begin{minipage}[t]{\codeW}
\begin{lstlisting}[style=fortran]
do step = 1, nsteps
   ireq        = 0
   recv_offset = 0
   do i = 0, size-1
      if (npoin_recv(i) > 0) then
         ireq = ireq + 1
         call MPI_Irecv( &
              recvbuf_all(recv_offset+1), &
              npoin_recv(i)*nfields, &
              MPI_DOUBLE_PRECISION, i, 0, &
              MPI_COMM_WORLD, &
              recv_requests(ireq), ierr)
         recv_offset = recv_offset + &
                       npoin_recv(i)*nfields
      end if
   end do
   call MPI_Waitall(nactive, recv_requests, &
                    recv_status, ierr)
   ! ... scatter into ordered_buf ...
end do
\end{lstlisting}
\end{minipage}\hfill
\begin{minipage}[t]{\codeW}
\begin{lstlisting}[style=julia]
# couplingStructs.jl :: coupling_exchange_data!
# Called once per Jexpresso time step from the
# DiscreteCallback in TimeIntegrators.jl.
function coupling_exchange_data!(cpg::CouplingData)
    send_requests = MPI.Request[]
    for dest_rank in cpg.send_to_ranks
        cpg.npoin_send[dest_rank+1] > 0 || continue
        push!(send_requests,
              MPI.Isend(cpg.send_bufs[dest_rank+1],
                        dest_rank, 0,
                        cpg.comm_world))
    end
    isempty(send_requests) ||
        MPI.Waitall(send_requests)
end
\end{lstlisting}
\end{minipage}

% ---------------------------------------------------------------- 8
\stepHdr{8.\ Receive-side scatter / interpolation pack}
\noindent
\begin{minipage}[t]{\stepW}
\footnotesize Alya re-orders the packed values into its flat-index slot using the global IDs saved at setup; Julia produces those values by SEM-interpolating onto the Alya points and packing per destination rank.
\end{minipage}\hfill
\begin{minipage}[t]{\codeW}
\begin{lstlisting}[style=fortran]
! Re-order the packed values from every Julia rank
! into Alya's flat-index slot using je_gids_all,
! the global IDs received once at setup.
ordered_buf = 0.0d0
do jpt = 1, total_pts
   local_pos = int(je_gids_all(jpt)) - 1 - i_start
   if (local_pos >= 0 .and. &
       local_pos < npoin_local) then
      ordered_buf(local_pos*nfields+1 : &
                  (local_pos+1)*nfields) = &
        recvbuf_all((jpt-1)*nfields+1 : jpt*nfields)
   end if
end do
\end{lstlisting}
\end{minipage}\hfill
\begin{minipage}[t]{\codeW}
\begin{lstlisting}[style=julia]
# couplingStructs.jl
# Per-step interpolation of the SEM solution to
# every local Alya grid point, then packing into
# per-destination send buffers in the same order
# the global-ID list was originally transmitted.
interpolate_solution_to_alya_coords!(
    u_interp, alya_coords, qout, xi_nodes, w, neqs,
    elem_bboxes, bins, e_conn, elem_x, elem_y,
    psi_xi, psi_eta, dpsi_xi, dpsi_eta,
    alpha, x_e, y_e, mesh_x, mesh_y)
pack_velocity_data!(cpg,
                    u_interp[:, 2:neqs-1],
                    owner_ranks)
\end{lstlisting}
\end{minipage}

\par\medskip
\noindent\textcolor{ruleC}{\rule{\linewidth}{0.3pt}}

\vspace{4pt}
{\footnotesize\itshape
Notes.  Both codes are launched as a single MPMD MPI job; Alya occupies
world ranks $0,\ldots,N_{\mathrm{alya}}-1$ and Jexpresso the remaining
ranks. All inter-code messages travel on \texttt{MPI\_COMM\_WORLD}; the
split communicators (\texttt{PAR\_COMM\_FINAL} on the Alya side,
\texttt{\_local\_comm} on the Julia side) are used only for intra-code
collectives. The collective order in steps 3--5 is identical on the two
sides --- this is what makes the coupling deterministic.\par}

\end{document}
